\documentclass[12pt]{article}
\usepackage[a4paper,margin=1in]{geometry}
\usepackage[utf8]{inputenc}
\usepackage[T1]{fontenc}
\usepackage{amsmath,amssymb}
\usepackage{textcomp}
\usepackage{parskip}
\begin{document}
\begin{center}
{\Large \textbf{2001 STEP II}}
\end{center}

\subsection*{Page 1}
\begin{flushleft}
STEP II, 2001 2\par
\bigskip\textbf{Section A: Pure Mathematics}\par
\bigskip\textbf{Question 1.} Use the binomial expansion to obtain a polynomial of degree 2 which is a good approximation\par
\[to \sqrt(1 x) when x is small.\]
-\par
\textbf{(i)} By taking x = 1/100, show that \sqrt(11) 79599/24000, and estimate, correct to 1\par
\approx\par
significant figure, the error in this approximation. (You may assume that the error\par
is given approximately by the first neglected term in the binomial expansion.)\par
\textbf{(ii)} Find a rational number which approximates \sqrt(1111) with an error of about 2 10\par
-12.\par
\[\times\]
\bigskip\textbf{Question 2.} Sketch the graph of the function [x/N ], for 0 < x < 2N , where the notation [y] means the\par
\[integer part of y. (Thus [2.9] = 2, [4] = 4.)\]
\textbf{(i)} Prove that\par
2N\par
\[( 1)[k/N]k = 2N N 2.\]
- -\par
\[Xk=1\]
\textbf{(ii)} Let\par
2N\par
\[S = ( 1)[k/N]2 -k.\]
N\par
-\par
\[Xk=1\]
Find S in terms of N and determine the limit of S as N .\par
N N\par
\to \infty\par
\bigskip\textbf{Question 3.} The cuboid ABCDEF GH is such AE, BF , CG, DH are perpendicular to the opposite faces\par
\[ABCD and EF GH, and AB = 2, BC = 1, AE = λ. Show that if α is the acute angle between\]
the diagonals AG and BH then\par
3 λ2\par
\[cos α = -\]
\{ 5 + λ2 \{\par
\{ \{\par
\{ \{\par
\{ \{\par
Let R be the ratio of the volume of the cuboid to its surface area. Show that R < 1/3 for all\par
possible values of λ.\par
\[Prove that, if R > 1/4, then α 6 arccos(1/9).\]
\end{flushleft}
\subsection*{Page 2}
\begin{flushleft}
STEP II, 2001 3\par
4 Let\par
\[f (x) = P sin x + Q sin 2x + R sin 3x .\]
\[Show that if Q^2 < 4R(P R), then the only values of x for which f (x) = 0 are given by\]
-\par
\[x = m\pi, where m is an integer.\]
\[[You may assume that sin 3x = sin x(4 cos2 x 1).]\]
-\par
Now let\par
\[g (x) = sin 2nx + sin 4nx sin 6nx,\]
-\par
where n is a positive integer and 0 < x < \pi/2. Find an expression for the largest root of the\par
\[equation g (x) = 0, distinguishing between the cases where n is even and n is odd.\]
\bigskip\textbf{Question 5.} The curve C passes through the origin in the x-y plane and its gradient is given by\par
\[dy = x(1 x^2)e -x^2 .\]
dx -\par
Show that C has a minimum point at the origin and a maximum point at 1, 1 e -1 . Find\par
1 2\par
the coordinates of the other stationary point. Give a rough sketch of C .\par
The curve C passes through the origin and its gradient is given by\par
\[dy = x(1 x^2)e -x^3 .\]
dx -\par
Show that C has a minimum point at the origin and a maximum point at (1, k), where\par
\[k > 1 e -1. (You need not find k.)\]
\bigskip\textbf{Question 6.} Show that\par
\[1 x^4 \pi 2\]
\[dx = .\]
\[Z 1 + x^2 4 - 3\]
Determine the values of\par
1 1 x\par
\textbf{(i)} x3 tan -1 - dx ,\par
Z \lambda 1 + x \mu\par
1 (1 y)3\par
\textbf{(ii)} - tan\par
-1y\par
dy .\par
Z (1 + y)5\par
\end{flushleft}
\subsection*{Page 3}
\begin{flushleft}
STEP II, 2001 4\par
\bigskip\textbf{Question 7.} In an Argand diagram, O is the origin and P is the point 2 + 0i. The points Q, R and S are\par
such that the lengths OP , P Q, QR and RS are all equal, and the angles OP Q, P QR and\par
QRS are all equal to 5\pi/6, so that the points O, P , Q, R and S are five vertices of a regular\par
12-sided polygon lying in the upper half of the Argand diagram. Show that Q is the point\par
\[2 + \sqrt\{3\} + i and find S.\]
The point C is the centre of the circle that passes through the points O, P and Q. Show\par
that, if the polygon is rotated anticlockwise about O until C first lies on the real axis, the\par
new position of S is\par
\[(3\sqrt\{2\} + \sqrt\{6\})(\sqrt\{3\} i) .\]
- 2 -\par
\bigskip\textbf{Question 8.} The function f satisfies f (x + 1) = f (x) and f (x) > 0 for all x.\par
\textbf{(i)} Give an example of such a function.\par
\textbf{(ii)} The function F satisfies\par
dF\par
\[= f (x)\]
dx\par
\[and F(0) = 0. Show that F(n) = nF(1), for any positive integer n.\]
\textbf{(iii)} Let y be the solution of the differential equation\par
dy\par
\[+ f (x)y = 0\]
dx\par
\[that satisfies y = 1 when x = 0. Show that y(n) 0 as n , where n = 1, 2, 3, . . .\]
\to \to \infty\par
\end{flushleft}
\subsection*{Page 4}
\begin{flushleft}
STEP II, 2001 5\par
\bigskip\textbf{Section B: Mechanics}\par
\bigskip\textbf{Question 9.} A particle of unit mass is projected vertically upwards with speed u. At height x, while the\par
particle is moving upwards, it is found to experience a total force F , due to gravity and air\par
\[resistance, given by F = α e\]
-βx,\par
where α and β are positive constants. Calculate the energy\par
expended in reaching this height. Show that\par
\[F = 1 βv2 + α 1 βu2 ,\]
2 - 2\par
where v is the speed of the particle, and explain why α = 1 βu2 + g, where g is the acceleration\par
due to gravity.\par
Determine an expression, in terms of y, g and β, for the air resistance experienced by the\par
particle on its downward journey when it is at a distance y below its highest point.\par
\bigskip\textbf{Question 10.} Two particles A and B of masses m and km, respectively, are at rest on a smooth horizontal\par
surface. The direction of the line passing through A and B is perpendicular to a vertical wall\par
which is on the other side of B from A. The particle A is now set in motion towards B with\par
speed u. The coefficient of restitution between A and B is e and between B and the wall is\par
e . Show that there will be a second collision between A and B provided\par
1 + e (1 + e )\par
2 1\par
\[k < .\]
e\par
Show that, if e = 1 , e = 1 and k < 5, then the kinetic energy of A and B immediately after\par
1 3 2 2\par
\[B rebounds from the wall is greater than mu2/27.\]
\bigskip\textbf{Question 11.} A two-stage missile is projected from a point A on the ground with horizontal and vertical\par
velocity components u and v, respectively. When it reaches the highest point of its trajectory\par
an internal explosion causes it to break up into two fragments. Immediately after this explosion\par
one of these fragments, P , begins to move vertically upwards with speed v , but retains the\par
e\par
previous horizontal velocity. Show that P will hit the ground at a distance R from A given\par
by\par
gR\par
\[= v + v + \sqrt(v^2 + v^2).\]
u e e\par
It is required that the range R should be greater than a certain distance D (where D > 2uv/g).\par
Show that this requirement is satisfied if\par
gD gD 2uv\par
\[v > - .\]
e\par
2u \lambda gD uv \mu\par
-\par
[The effect of air resistance is to be neglected.]\par
\end{flushleft}
\subsection*{Page 5}
\begin{flushleft}
STEP II, 2001 6\par
\bigskip\textbf{Section C: Probability and Statistics}\par
\bigskip\textbf{Question 12.} The national lottery of Ruritania is based on the positive integers from 1 to N , where N is\par
very large and fixed. Tickets cost \pounds 1 each. For each ticket purchased, the punter (i.e. the\par
purchaser) chooses a number from 1 to N . The winning number is chosen at random, and the\par
jackpot is shared equally amongst those punters who chose the winning number.\par
A syndicate decides to buy N tickets, choosing every number once to be sure of winning\par
a share of the jackpot. The total number of tickets purchased in this draw is 3.8N and the\par
jackpot is \pounds W . Assuming that the non-syndicate punters choose their numbers independently\par
and at random, find the most probable number of winning tickets and show that the expected\par
net loss of the syndicate is approximately\par
5(1 e\par
-2.8)\par
N - W .\par
- 14\par
\bigskip\textbf{Question 13.} The life times of a large batch of electric light bulbs are independently and identically dis-\par
tributed. The probability that the life time, T hours, of a given light bulb is greater than t\par
hours is given by\par
\[P(T > t) = ,\]
(1 + kt)α\par
\[where α and k are constants, and α > 1. Find the median M and the mean m of T in terms\]
of α and k.\par
Nine randomly selected bulbs are switched on simultaneously and are left until all have failed.\par
The fifth failure occurs at 1000 hours and the mean life time of all the bulbs is found to be 2400\par
hours. Show that α 2 and find the approximate value of k. Hence estimate the probability\par
\approx\par
that, if a randomly selected bulb is found to last M hours, it will last a further m M hours.\par
-\par
\end{flushleft}
\subsection*{Page 6}
\begin{flushleft}
STEP II, 2001 7\par
\bigskip\textbf{Question 14.} Two coins A and B are tossed together. A has probability p of showing a head, and B has\par
\[probability 2p, independent of A, of showing a head, where 0 < p < 1 . The random variable\]
X takes the value 1 if A shows a head and it takes the value 0 if A shows a tail. The random\par
variable Y takes the value 1 if B shows a head and it takes the value 0 if B shows a tail. The\par
random variable T is defined by\par
\[T = λX + 1 (1 λ)Y.\]
2 -\par
Show that E(T ) = p and find an expression for Var (T ) in terms of p and λ. Show that as λ\par
varies, the minimum of Var (T ) occurs when\par
1 2p\par
\[λ = - .\]
3 4p\par
-\par
The two coins are tossed n times, where n > 30, and T is the mean value of T . Let b be a fixed\par
\[positive number. Show that the maximum value of P( T p < b) as λ varies is approximately\]
| - |\par
2Φ(b/s) 1, where Φ is the cumulative distribution function of a standard normal variate and\par
-\par
p(1 p)(1 2p)\par
\[s^2 = - - .\]
(3 4p)n\par
-\par
\end{flushleft}
\end{document}
